\begin{table}[H]
\caption{Сравнение алгоритмов поиска}
\label{tbl:search-comparison}
\centering
\setlength{\tabcolsep}{4pt}
\begin{NiceTabular}{C{0.21}C{0.15}C{0.15}C{0.24}C{0.25}}[hvlines]
\Block{2-1}{\diagbox{\bfseries Алгоритм}{\bfseries Критерии}}
  & \Block{1-2}{\bfseries Оценка временной\\ \bfseries эффективности как\\ \bfseries функции от числа\\ \bfseries необходимых\\ \bfseries сравнений}
  &
  & \Block{2-1}{\bfseries Требования,\\ \bfseries налагаемые\\ \bfseries алгоритмом}
  & \Block{2-1}{\bfseries Оценка памяти,\\ \bfseries затрачиваемой\\ \bfseries алгоритмом сверх\\ \bfseries входных данных} \\
  & \textbf{ЛС} & \textbf{ХС} & & \\
Поиск полным перебором
  & $1$
  & $N$
  & Нет
  & $i \to O(1)$, 1 вызов функции \\
Нерекурсивный бинарный поиск
  & $1 + \log_2 N$
  & $1 + \log_2 N$
  & \Block{3-1}{Массив\\ предварительно\\ отсортирован\\ и поддерживается\\ отсортированным}
  & \Block{2-1}{$l, r, m \to O(1)$,\\ 1 вызов функции} \\
Нерекурсивный бинарный поиск (предварительный выход)
  & $1$
  & $2 \log_2 N$
  & & \\
Рекурсивный бинарный поиск
  & $1$
  & $2 \log_2 N$
  &
  & $\log_2 N$ вызовов, т.\,е. $\log_2 N \cdot {}$(размер кадра стека) памяти \\
\end{NiceTabular}
\end{table}